% U5 WS1 -- rates and rate laws. Blocks 1-2.
\documentclass{apchem-worksheet}

\apcunit{5}
\apcunittitle{Kinetics}
\apcdoctitle{Worksheet 1 \textbullet\ Rates and Rate Laws}
\apcsubtitle{CED 5.1--5.2 \textbullet\ orders come from data, never from coefficients}

\begin{document}
\wsheader[{$\text{rate} = -\tfrac{1}{a}\Delta[\ce{A}]/\Delta t =
+\tfrac{1}{c}\Delta[\ce{C}]/\Delta t$} \textbullet\
{$\text{rate} = k[\ce{A}]^m[\ce{B}]^n$} \textbullet\ orders are determined
experimentally.]

\problem[]{\ced{5.1} For $\ce{2N2O5(g) -> 4NO2(g) + O2(g)}$, oxygen forms at
\SI{2.0e-3}{\Molar\per\second}.}
\begin{parts}
  \item The rate of the reaction: \answerblank[3.4cm]{\SI{2.0e-3}{\Molar\per\second}}
  \item Rate of disappearance of \ce{N2O5}:
        \answerblank[3.4cm]{\SI{4.0e-3}{\Molar\per\second}}
  \item Rate of appearance of \ce{NO2}:
        \answerblank[3.4cm]{\SI{8.0e-3}{\Molar\per\second}}
  \item Why do the three numbers differ? \answerlines[2]{Each species is
        consumed or formed at a rate set by its coefficient. Dividing each
        by its coefficient gives one common value --- the reaction rate ---
        which is why the convention exists.}
\end{parts}

\problem[]{\ced{5.1} For $\ce{N2 + 3H2 -> 2NH3}$, \ce{H2} is consumed at
\SI{0.030}{\Molar\per\second}. Find the reaction rate and the rate of
formation of \ce{NH3}. \workspace[2.4cm]
\keyonly{\begin{apcanswer}rate $= (1/3)(0.030) =
\mathbf{0.010}~\si{\Molar\per\second}$;
\ce{NH3} forms at $2(0.010) =
\mathbf{0.020}~\si{\Molar\per\second}$\end{apcanswer}}}

\problem[]{\ced{5.1} A plot of $[\ce{A}]$ against time is steep at first and
flattens.}
\begin{parts}
  \item Average rate is the slope of: \answerblank[8.7cm]{a straight line
        between two points}
  \item Instantaneous rate is the slope of:
        \answerblank[5.2cm]{the tangent at that point}
  \item Why does the curve flatten? \answerlines[2]{Reactant is being
        consumed, so its concentration --- and therefore the frequency of
        reactive collisions --- decreases with time.}
  \item Why are rate laws determined from \emph{initial} rates?
        \answerlines[2]{At $t=0$ no product has built up, so there is no
        reverse reaction and no interference from products to complicate the
        measurement.}
\end{parts}

\problem[]{\ced{5.2} Use the data to find the rate law and $k$.}
\begin{center}
\begin{tabular}{@{}cccc@{}}
\hline
\textbf{Exp} & $[\ce{A}]_0$ (M) & $[\ce{B}]_0$ (M) &
  \textbf{rate} (\si{\Molar\per\second}) \\
\hline
1 & 0.10 & 0.10 & $2.0\times10^{-3}$ \\
2 & 0.20 & 0.10 & $8.0\times10^{-3}$ \\
3 & 0.10 & 0.20 & $4.0\times10^{-3}$ \\
\hline
\end{tabular}
\end{center}
\begin{parts}
  \item Order in A, with the comparison you used: \answerlines[2]{2.
        Comparing experiments 1 and 2, only $[\ce{A}]$ changed: doubling it
        multiplied the rate by 4, and $4 = 2^2$.}
  \item Order in B: \answerblank[1.8cm]{1}
  \item Overall order: \answerblank[1.8cm]{3}
  \item Rate law: \answerblank[4.4cm]{rate $= k[\ce{A}]^2[\ce{B}]$}
  \item Calculate $k$ with units. \workspace[2cm]
        \keyonly{\begin{apcanswer}$k = 2.0\times10^{-3}/[(0.10)^2(0.10)] =
        2.0\times10^{-3}/1.0\times10^{-3} =
        \mathbf{2.0}~\si{\per\Molar\squared\per\second}$\end{apcanswer}}
\end{parts}

\problem[]{\ced{5.2} A second reaction:}
\begin{center}
\begin{tabular}{@{}cccc@{}}
\hline
\textbf{Exp} & $[\ce{X}]_0$ (M) & $[\ce{Y}]_0$ (M) &
  \textbf{rate} (\si{\Molar\per\second}) \\
\hline
1 & 0.20 & 0.10 & $1.5\times10^{-2}$ \\
2 & 0.40 & 0.10 & $3.0\times10^{-2}$ \\
3 & 0.20 & 0.20 & $1.5\times10^{-2}$ \\
\hline
\end{tabular}
\end{center}
\begin{parts}
  \item Rate law: \answerblank[4cm]{rate $= k[\ce{X}]$}
  \item $k$ with units: \answerblank[3cm]{\SI{0.075}{\per\second}}
  \item Y is a reactant yet does not appear in the rate law. Explain what
        that tells you. \answerlines[3]{Y is zero order --- changing its
        concentration does not change the rate. Mechanistically, Y does not
        take part in the rate-determining step, even though it is consumed
        somewhere in the overall reaction.}
\end{parts}

\problem[]{\ced{5.2} Units as a check. Give the units of $k$ for:}
\begin{parts}
  \item zero order: \answerblank[3cm]{\si{\Molar\per\second}}
  \item first order: \answerblank[3cm]{\si{\per\second}}
  \item second order: \answerblank[3cm]{\si{\per\Molar\per\second}}
  \item Why are the units a useful self-check?
        \answerlines[2]{Rate always has units of \si{\Molar\per\second}, so
        the units of $k$ are fixed by the overall order. If your $k$ carries
        the wrong units, your orders must be wrong.}
\end{parts}

\problem[]{\ced{5.2} For rate $= k[\ce{A}]^2[\ce{B}]$ with
$k = \SI{2.0}{\per\Molar\squared\per\second}$:}
\begin{parts}
  \item Find the rate when $[\ce{A}] = 0.30$ and $[\ce{B}] = 0.20$.
        \workspace[1.8cm]
        \keyonly{\begin{apcanswer}$(2.0)(0.30)^2(0.20) =
        \mathbf{0.036}~\si{\Molar\per\second}$\end{apcanswer}}
  \item Predict the factor by which the rate changes if $[\ce{A}]$ is
        tripled and $[\ce{B}]$ halved.
        \answerblank[3cm]{$\times 4.5$}
\end{parts}

\problem[]{\ced{5.2} Zumdahl's decomposition $\ce{2N2O5 -> 4NO2 + O2}$ is
experimentally \emph{first} order in \ce{N2O5}. Explain why this single fact
disposes of the idea that coefficients give orders.
\answerlines[3]{If coefficients gave orders, this reaction would be second
order in \ce{N2O5}. It is measured to be first order. Since the overall
equation only summarizes stoichiometry while the rate law reflects the
slowest elementary step, there is no reason for the two to agree --- and
here they demonstrably do not.}}

\begin{spiralstrip}{Units 1--4 \textbullet\ spiral}
  \item Net ionic for \ce{Ag+} with \ce{Cl-}:
        \answerblank[5.0cm]{\ce{Ag+ + Cl- -> AgCl(s)}}
  \item Oxidation state of S in \ce{SO4^2-}: \answerblank[1.8cm]{$+6$}
  \item Strongest IMF in \ce{H2O}: \answerblank[3.8cm]{hydrogen bonding}
  \item Moles in \SI{25.0}{\milli\litre} of \SI{0.200}{\Molar}:
        \answerblank[2.6cm]{\SI{0.00500}{\mole}}
  \item Electron configuration of \ce{Ca}:
        \answerblank[3.4cm]{[Ar]$4s^2$}
\end{spiralstrip}

\end{document}
